%! AP Statistics — Unit 9, Lesson 9.2 — Group Activity (Answer Key)
\documentclass[10pt]{article}
\usepackage{apstats-article}
\usepackage{apstats-key}

\begin{document}

\pageheader{Unit 9, Lesson 9.2}{Group Activity --- Answer Key}
\namepartnerperiod

\vspace{0.1in}
\textbf{Building a CI for the Slope --- Car Weight and Fuel Efficiency}

\vspace{0.05in}
A researcher randomly selects 30 cars from a large database and records curb weight
($x$, in hundreds of pounds) and highway fuel efficiency ($y$, mpg).
Computer output is shown below.

\vspace{0.1in}
\begin{center}
\renewcommand{\arraystretch}{1.3}
\begin{tabular}{lrrrr}
    \textbf{Predictor} & \textbf{Coef} & \textbf{SE Coef} & \textbf{T} & \textbf{P} \\
    \hline
    Constant & 47.8 & 3.4 & 14.06 & 0.000 \\
    Weight   & $-2.3$ & 0.41 & $-5.61$ & 0.000 \\
    \hline
    \multicolumn{5}{l}{$s = 3.2$ \quad $R^2 = 53.2\%$ \quad $n = 30$}
\end{tabular}
\end{center}

\vspace{0.05in}
The researcher wants to build a 95\% CI for the true slope $\beta$.
Use $t^* = 2.048$ ($df = 28$).

\vspace{0.1in}

\begin{tcolorbox}[colback=white, colframe=black!40,
    title=\textbf{Tier R --- Remediate},
    fonttitle=\bfseries, arc=2mm, left=3mm, right=3mm, top=2mm, bottom=2mm]

\begin{enumerate}[label=\alph*.]

    \item Identify the slope estimate and standard error from the output.

          $b = $ \ans{$-2.3$}\quad and\quad $SE_b = $ \ans{$0.41$}

    \item Complete the confidence interval.

          $b \pm t^* \cdot SE_b
          = $ \ans{$-2.3$} $\pm\ 2.048 \times $ \ans{$0.41$}
          $= -2.3 \pm$ \ans{$0.840$}
          $= ($ \ans{$-3.14$},\; \ans{$-1.46$}$)$

    \item Fill in the interpretation.

          ``We are 95\% confident that the true slope relating
          \ans{highway fuel efficiency (mpg)} to \ans{curb weight (hundreds of pounds)}
          for \ans{all cars in the database}
          is between \ans{$-3.14$} and \ans{$-1.46$}
          mpg per hundred pounds.''

\end{enumerate}
\end{tcolorbox}

\begin{tcolorbox}[colback=white, colframe=black!40,
    title=\textbf{Tier A --- Approaching Proficiency},
    fonttitle=\bfseries, arc=2mm, left=3mm, right=3mm, top=2mm, bottom=2mm]

\begin{enumerate}[label=\alph*.]

    \item \ans{$b = -2.3$ mpg per hundred pounds; $SE_b = 0.41$.}

    \item \ans{\textbf{L}: Residual plot shows random scatter with no curve --- linear condition met.
          \textbf{I}: Random sample from a large database; $n = 30 \le 10\%$ of all cars --- independent.
          \textbf{N}: $n = 30$, which is large enough to invoke the Central Limit Theorem --- normal condition met.
          \textbf{E}: Residual plot shows roughly constant vertical spread --- equal variance met.
          \textbf{R}: Cars were randomly selected --- random condition met.}

    \item \ans{$CI: -2.3 \pm 2.048(0.41) = -2.3 \pm 0.840 = (-3.14,\ -1.46)$ mpg per hundred pounds.}

    \item \ans{We are 95\% confident that the true slope of the population regression
          line relating highway fuel efficiency to curb weight for all cars in the database
          is between $-3.14$ and $-1.46$ mpg per hundred pounds.}

    \item \ans{Yes. Since 0 is not contained in the interval $(-3.14,\ -1.46)$, this provides
          evidence that $\beta \ne 0$; there is a negative linear relationship between
          car weight and fuel efficiency.}

\end{enumerate}
\end{tcolorbox}

\begin{tcolorbox}[colback=white, colframe=black!40,
    title=\textbf{Tier E --- Extension},
    fonttitle=\bfseries, arc=2mm, left=3mm, right=3mm, top=2mm, bottom=2mm]

Complete all of Tier A above. Then answer the following.

\begin{enumerate}[label=\alph*., start=6]

    \item \ans{The CI would be wider. The formula $SE_b = s/(s_x\sqrt{n-1})$ shows that a
          smaller $n$ produces a larger $SE_b$, which increases the margin of error and widens
          the interval. With $n = 15$, we also lose degrees of freedom, making $t^*$ larger,
          which further widens the CI.}

    \item \ans{Yes, $-2.0$ is inside the interval $(-3.14,\ -1.46)$, so the manufacturer's
          claimed value of $\beta = -2.0$ is consistent with the data. The CI does not
          provide evidence against their claim.}

\end{enumerate}
\end{tcolorbox}

\end{document}
